\chapter{Cable modelling}
\label{cable_modelling}

\minitoc{}

\section{Introduction}

This chapter focuses on the modeling of a cable --particularly its shape-- attached to underwater
robots under different usage modes. With an accurate model, the cable can be usued as a
localization method for the robots as well as a basis for control constraints. Its usage in the
control scheme of the system constrains the model to be as computationally efficient as possible
while staying accurate. The model must also be inferred from available sensors: IMUs, water
pressure sensor and cameras.

The contributions of this chapter are the following:
\begin{itemize}
	\item
\end{itemize}

\section{Catenary model}

\subsection{Definition}

\lipsum[1]

\subsection{Parameter derivation}

\lipsum[1]

\section{Augmented catenary model}

\subsection{Definition}

\lipsum[1]

\subsection{Parameters derivation}

\lipsum[1]

\subsection{Parameters dynamics}

\lipsum[1]

\section{Conclusion}

\lipsum[1]